\documentclass[12pt,a4paper]{exam}
\usepackage[utf8]{inputenc}
\usepackage[T1]{fontenc}
\usepackage{amsmath}
\usepackage{amsfonts}
%\usepackage{amssymb}
\usepackage{graphicx}
\usepackage{geometry}
\usepackage{enumitem}

\geometry{a4paper, margin=2cm}

\usepackage{cprotect}

\usepackage{xcolor}
\definecolor{maroon}{cmyk}{0, 0.87, 0.68, 0.32}
\definecolor{halfgray}{gray}{0.55}
\definecolor{ipython-frame}{RGB}{207, 207, 207}
\definecolor{ipython-bg}{RGB}{247, 247, 247}
\definecolor{ipython-red}{RGB}{186, 33, 33}
\definecolor{ipython-green}{RGB}{0, 128, 0}
\definecolor{ipython-cyan}{RGB}{64, 128, 128}
\definecolor{ipython-purple}{RGB}{170, 34, 255}

\usepackage{listings}
\lstdefinelanguage{iPython}{
	morekeywords={access,and,del,except,exec,in,is,lambda,not,or,raise},
	morekeywords=[2]{for,print,abs,all,any,basestring,bin,bool,bytearray,callable,chr,classmethod,cmp,compile,complex,delattr,dict,dir,divmod,enumerate,eval,execfile,file,filter,float,format,frozenset,getattr,globals,hasattr,hash,help,hex,id,input,int,isinstance,issubclass,iter,len,list,locals,long,map,max,memoryview,min,next,object,oct,open,ord,pow,property,range,reduce,reload,repr,reversed,round,set,setattr,slice,sorted,staticmethod,str,sum,super,tuple,type,unichr,unicode,vars,xrange,zip,apply,buffer,coerce,intern,elif,else,if,continue,break,while,class,def,return,try,except,import,finally,try,except,from,global,pass, True, False},
	sensitive=true,
	morecomment=[l]\#,%
	morestring=[b]',%
	morestring=[b]",%
	moredelim=**[is][\color{black}]{@@}{@@},
	identifierstyle=\color{black}\footnotesize\ttfamily,
	commentstyle=\color{ipython-cyan}\footnotesize\itshape\ttfamily,
	stringstyle=\color{ipython-red}\footnotesize\ttfamily,
	keepspaces=true,
	showspaces=false,
	showstringspaces=false,
	rulecolor=\color{ipython-frame},
	frame=single,
	frameround={t}{t}{t}{t},
	backgroundcolor=\color{ipython-bg},
	basicstyle=\footnotesize\ttfamily,
	keywordstyle=[2]\color{ipython-green}\bfseries\footnotesize\ttfamily, 
	keywordstyle=\color{ipython-purple}\bfseries\footnotesize\ttfamily
}

\lstdefinelanguage{iOutput} {
	sensitive=true,
	identifierstyle=\color{black}\small\ttfamily,
	stringstyle=\color{ipython-red}\small\ttfamily,
	keepspaces=true,
	showspaces=false,
	showstringspaces=false,
	rulecolor=\color{ipython-frame},
	basicstyle=\small\ttfamily,
}

\lstnewenvironment{ipython}[1][]{\lstset{language=iPython,mathescape=true,escapeinside={*@}{@*}}%
}{%
}

\lstnewenvironment{ioutput}[1][]{\lstset{language=iOutput,mathescape=true,escapeinside={*@}{@*}}%
}{%
}


\title{Financial Market Course 25/26\\ Exam}
\author{Prof. Simone Freschi, Prof. Matteo Sani}
\date{$09^{\mathrm{st}}$ January 2026}

\printanswers
%\noprintanswers
\begin{document}
\maketitle

\begin{center}
\fbox{\fbox{\parbox{5.5in}{
Answer the questions in the provided spaces. If you run out of room for an answer, continue on the page back.}}}
\end{center}

\begin{center}
\vspace{5mm}
\makebox[0.75\textwidth]{Student's name:\enspace\hrulefill}
\end{center}

\section*{Questions}
\vspace{.5cm}
\begin{questions}

\question 
Credit Risk and CVA.
\begin{enumerate}
\item Describe the probability of the default formula using the intensity model. The counterparty BB Company LTD has a CDS equal to 4\%; calculate the probability of default at 1y horizon using the shortcut formula for the intensity value (Recovery Value equal to 40\%).	
\item Describe the CVA adjustment for interest rate swaps. Imagine that BB Company LTD enters into a EUR 100mln 5 years swap (where it receives fixed 5\%), and subsequently, rates increases to 6\%; what will happen to the mark to market of the swap (if you can calculate the approximante change in fair value) ? What will happen to the CVA all the rest being equal ? In case of a decrease of the CDS to 2\% do you expect a different answer ?\end{enumerate}
\fillwithlines{3cm}
\begin{solution}
\end{solution}

\question Describe the asset swap contract for a coupon bond with coupons equal to $C$ and dirty Price equal to $P(t)$.
\begin{enumerate}
\item What will happen to the price if the sprea increase, all the rest being equal ? What does it mean in term of credit quality expectation by the market for the bond's issuer ?
\item What is the difference between asset swap spread and \textbf{G-spread} or \textbf{I-spread} calculated in Bloomberg screen ?
\end{enumerate}
\fillwithlines{3cm}
\begin{solution}
\end{solution}

\question Describe the equation for the return atribution model by Brinson, Hood and Beebower (BHB).
\begin{enumerate}[label=(\alph*)]
\item Explain the difference between Allocation Effect and Security Selection Effect;
\item Imagine that your portfolio is overweight Energy (30\% vs 20\% Benchmark) and the Energy sector is up 15\%. Your energy stock in the portfolio are up only 5\%. Determine the allocation effect and the selection effect and the overall portfolio result.
\end{enumerate}
\fillwithlines{3cm}
\begin{solution}
\end{solution}

\question Risk-Performance Evaluation Measures.
\begin{enumerate}
\item Describe the Capture Ratio Measures. Imagine that Asset manager \textbf{A} has an Upside Capture Ratio of 140\% and Downside Capture Ratio of 110\%; Asset manager \textbf{B} has an Upside Capture Ratio of 90\% and Downside Capture Ratio of 70\%; who is the best according to this measure?
\item Describe the Draw-Down measure and the Max Draw-Down measure. Imagine that Asset manager \textbf{A} has an expected excess return of 10\% and Max Draw-Down equal to 20\%; Asset manager \textbf{B} has an expected excess return of 5\% and Max Draw-Down equal to 7\%; who would you prefer and why? Will your choice be different if you had a target return of more or equal to 10\%?
\end{enumerate}
\fillwithlines{3cm}


\question
Consider two \textbf{independent} assets, A and B. Both are "digital" assets with the following loss profiles $L$ over 1 year horizon:
\begin{itemize}
\item A: has a 4\% probability of a loss of \$100 million. Otherwise, the loss is \$0;
\item B: has a 2\% probability of a loss of \$500 million. Otherwise, the loss is \$0.
\end{itemize}

Calculate:
\begin{parts}
\part 95\% c.l. VaR for each asset A and B;
\part VaR$_{95\%}$ for a portfolio $P=A+B$. Check if the sub-additivity property is fullfilled;
\part 95\% c.l. Expected Shortfall (ES) for each asset separately and portfolio $P$. Compare these results with those of parts (a) and (b).
\end{parts}

%\fillwithdottedlines{5cm}
\makeemptybox{15cm}
\begin{solution}
\begin{parts}
\part 95\% c.l. VaR is the 95th quantile of the loss distribution. Hence for both the two assets is \$0 since 95\% of the distribution is in the 0 loss "region" (0.95 < 0.96).

VaR only tells us the "threshold", a risk manager looking only at VaR would see these portfolios as identical in risk.

\part For portfolio $P$ we have the following loss distribution (the assets are independent):
\begin{itemize}
\item \$0 loss: $0.96 \cdot 0.96 = 0.9216$;
\item \$300 M loss: $2 \cdot (0.04 \cdot 0.96) = 0.0768$;
\item \$600 M loss: $0.04 \cdot 0.04 = 0.0016$.
\end{itemize}

\begin{equation*}
\text{VaR}_{95\%}(P)=300M \ge \text{VaR}_{95\%}(A) + \text{VaR}_{95\%}(B) = 0M
\end{equation*}

VaR violates sub-additivity. The risk of losses increases diversifying the portfolio, which is in contrast with basic financial principles.

\part ES is the average loss in the worst $(1-\alpha)\%$ quantile. For the single assets we get

\begin{gather*}
	\text{ES}(A) = \frac{0.04\cdot 100 + 0.01\cdot 0}{0.05} = \$80 M \\
	\text{ES}(B) = \frac{0.04\cdot 500 + 0.01\cdot 0}{0.05} = \$400 M \\	
\end{gather*}

ES looks at the entire area under the tail. It correctly identifies Portfolio B as significantly riskier because it accounts for the magnitude of the "blow-up" beyond the threshold.

For the portfolio $P$ instead:
\begin{equation*}
	\text{ES}(P) = \frac{0.0484\cdot 300 + 0.0016\cdot 600}{0.05} = \$309 M 
\end{equation*}

Note that ES is mathematically proven to be a Coherent Risk Measure. It always satisfies sub-additivity, ensuring that a diversified portfolio never appears riskier than the sum of its parts.
\end{parts}

\end{solution}

%%%%%%%%%%%%%%%%%%%%%%%%%%%%%%%%%%%%%
\question 
Bank \textit{BestBankEver} owns a portfolio containing $n=5$ loans. The portfolio is split into two classes:

\begin{itemize}
	\item Low risk: 2 loans with $p_{\text{def}}^L = 0.01$;
	\item High risk: 3 loans with $p_{\text{def}}^H = 0.05$.
\end{itemize}

\begin{parts}
\part Calculate the probability that all loans default simultaneously assuming independency ($\rho = 0$);
\part calculate the probability that all loans default simultaneously assuming perfect correlation ($\rho = 1.0$);
\part imagine then that the bank decides to make the portfolio safer by adding 5 more "Low risk" loans (i.e. we now have $n=10$ loans: 7 Low risk, 3 High risk).
In case of a scenario, like 2008 housing bubble burst, where $\rho$ tends to 1 how much the risk (default probability) is reduced ?
\end{parts}

\makeemptybox{5cm}
\fillwithdottedlines{5cm}
\begin{solution}
\begin{parts}
	\part In case of independent assets the default probability is just the product of each marginal probability:
		\begin{equation*}
			P_{\text{def}} = 0.01^2\cdot 0.05^3 = 1.25\cdot 10^{-8}
		\end{equation*}  
		which represents a super-rare event.
	\part If there is perfect correlation all assets defaults at once, but since there are two levels of probability this global event is driven by the stronger assets (low risk ones):
		\begin{equation*}
			P_{\text{def}} = \min(0.01, 0.05) = 0.01
		\end{equation*}  
		Note how, despite the same risk profile, the correlation alone increases the default probability by 6 order of magnitude.
	\part By increasing the number of \textit{safer} assets the bank is not actually reducing the risk in the particular bad scenario with $\rho$ close to 1.
	Indeed following the reasoning of part (b) the default probability is still given by $P_{\text{def}} = \min(P_{low}, P_{high}) = 0.01$ so the action taken by \textit{BestBankEver}
	just increases the number of assets simultaneously defaulting.
\end{parts}
\end{solution}

%%%%%%%%%%%%%%%%%%%%%%%%%%%%%%%%%%%%%
\question
A quant is using a Monte Carlo simulation to estimate the probability of two events occurring simultaneously. Each event as a marginal probability of $p=5\%$ and they are independent.

\begin{parts}
\part The analyst decides to run a Monte Carlo using \textbf{10 trials}. She uses a random number generator that yields the following pairs of $\mathcal{U}(0,1)$ variables for event A and B:
\begin{align*}
	(0.12,0.45),(0.88,0.21),(0.04,0.03),(0.55,0.91),(0.32,0.11),\\
	(0.67,0.44),(0.09,0.52),(0.22,0.18),(0.95,0.82),(0.41,0.33)
\end{align*}

How many "successes" occur in these 10 trials ? What is the relative error with respect to the \textit{true} probability for the events to happen ?

\part Which action should the manager take if she needs to reduce to a 10\% relative uncertainty ?
\end{parts}
\textbf{Hint}: for a proportion (Bernoulli distribution) the standard deviation $\sigma$ is $\sqrt{p(1-p)}$.
\makeemptybox{10cm}

\begin{solution}
\begin{parts}
\part Given the program output the estimated probability is $1/10=10\%$. This is double the individual $p$ and vastly different from the theoretical 
probability $p^2=0.25\%$. Low-frequency events (tails) require massive sample sizes. Small MC simulations lead to "sampling error" where the tail is either 
completely missed or wildly overestimated.

The relative error turns out to be huge:
\begin{equation*}
p_{true} = \frac{|p_{sim}-p_{true}|}{p_{true}}=\frac{|0.10-0.0025|}{0.0025}=3900 \%
\end{equation*}
\part Start by setting $\cfrac{\epsilon}{p_{true}} = 0.1$, then
\begin{align}
	\cfrac{\sqrt{\frac{p_{true}(1-p_{true})}{n}}}{p_{true}} = 0.10 \\
	\sqrt{\cfrac{1-p_{true}}{np_{true}}} = 0.10 \\
	\cfrac{1-p_{true}}{np_{true}} = 0.01 \\
	n = \cfrac{1-0.0025}{0.01\cdot 0.0025} \approx 40000
\end{align}
This demonstrates that estimating rare events (tails) requires a massive large amount  of computational power to be reliable.
\end{parts}	
\end{solution}

%%%%%%%%%%%%%%%%%%%%%%%%%%%%%%%%%%%%%

\question
You have been given the following liquid instruments in the market. All bonds pay coupons annually, and the face value is \$100.

\begin{table}[h]
	\centering
	\begin{tabular}{@{}llll@{}}
		\hline
		\textbf{Instrument} & \textbf{Maturity} & \textbf{Coupon Rate} & \textbf{Market Price} \\ 
		\hline
		Bond A (Zero)   & 1 Year & 0\% & \$97.00 \\
		Bond B (Coupon) & 2 Years & 5\% & \$101.00 \\
		Bond C (Coupon) & 3 Years & 6\% & \$102.00 \\ 
		\hline
	\end{tabular}
\end{table}

\begin{parts}
\part Using the bootstrap technique calculate the spot rates for years 1, 2, and 3 $(r_1, r_2, r_3)$.
\part Suppose there was a data entry error. The price of Bond A was actually \$96.50 instead of \$97.00 (a small 0.5\% difference). Without recalculating everything, describe what happens to $r_3$ if $r_1$ is calculated incorrectly.
\part You have the spot rates for Year 2 and Year 3. A manager wants to value a payment occurring at Year 2.5 (he suggests using a simple Linear Interpolation between $r_2$ and $r_3$). Although the spot rate curve might seem smooth what happens to the 6m-forward rate curve at 2.5 years ?
\end{parts}

\makeemptybox{5cm}
\fillwithdottedlines{5cm}

\begin{solution}	
\begin{parts}
\part Year 1 Spot Rate ($r_1$)
\begin{equation*}
	97=\frac{100}{(1+r_1)^1} \implies 1+r_1 = \frac{100}{97} \implies r_1 \approx 3.09\%
\end{equation*}
	Year 2 Spot Rate ($r_2$), using the result from $r_1$ to value Bond B:		
	\begin{equation*}
		101=\frac{5}{1.03093}+\frac{105}{(1+r_2)^2} \implies r_2 = 4.50\%
	\end{equation*}
	Year 3 Spot Rate ($r_3$):
	\begin{equation*}
		102=\frac{6}{1+r_1}+\frac{6}{(1+r_2)^2}+\frac{106}{(1+r_3)^3} \implies r_3 \approx 5.34\%		
	\end{equation*}
\part Because bootstrapping is a recursive process, an error in the Year 1 rate "contaminates" the Year 2 rate, which then compounds into the Year 3 rate. A small mispricing in short-term liquid instruments can lead to wildly unrealistic long-term forward rates.
\part 
\begin{equation*}
	r_{2.5} = \frac{(5.34-4.50)}{2} = 4.92 \%
\end{equation*}
Consequently:
\begin{gather*}
	f_{2,2.5} = \frac{2.5*4.92 - 2*4.50}{0.5} = 6.6 \% \\
	f_{2.5,3} = \frac{3*5.34 - 2.5*4.92}{0.5} = 7.44 \% \\
\end{gather*}
So while the spot curve might look smooth with linear interpolation, the resulting 6m-forward curve  will have "kinks" or "steps." In finance, we assume rates change continuously; a "steppy" forward curve creates arbitrage opportunities in the model that don't exist in reality.
\end{parts}
\end{solution}

%%%%%%%%%%%%%%%%%%%%%%%%%%%%%%%%%%%%%
\question 
At $t=0$ a bank enters into a 1-year Forward Contract to buy 1,000 barrels of oil from a counterparty at a delivery price of $K=\$80$ per barrel.
The counterparty is a risky corporate entity with a constant annual probability of default (PD) of 5\%.
The recovery rate $R=40\%$ and the interest rate $r$ is very close to 0\%.

The bank models the oil price at the end of the year ($t=1$) using only three equally likely scenarios:
\begin{itemize}
\item scenario A (Bull): $S_1=\$110$;
\item scenario B (Base): $S_2=\$80$;
\item scenario C (Bear): $S_3=\$50$.
\end{itemize}

\begin{parts}
\part Calculate the "Expected Exposure" (EE) of the bank at maturity and the CVA for this 1-year contract.
\part The \textit{Wrong-Way Risk} (WWR) occurs when the exposure to a counterparty is positively correlated with that counterparty's probability of default.
Suppose the counterparty is an oil producer, and let's drop the assumption of a PD independent of the oil price. Is the bank affected by WWR ?
\part If the oil price goes to \$50, the bank owes the counterparty money. What is the name of the adjustment the counterparty should make?
\end{parts}
%\textbf{Hint}: the bank only faces credit risk if the contract has a positive value.
\makeemptybox{10cm}

\begin{solution}
\begin{parts}
\part Calculate the mark-to-market for each scenario:
\begin{itemize}
\item A: $1,000\cdot(110-80)=+\$30,000$;
\item B: $1,000\cdot (80-80)=\$0$;
\item C: $1,000\cdot (50-80)=-\$30,000$.
\end{itemize}
then identify the exposure ($E=\max(MtM,0)$):
\begin{itemize}
\item A: \$30,000;
\item B: \$0;
\item C: \$0.
\end{itemize}
Finally calculate $EE$:
\begin{equation*}
EE=\frac{1}{3}(30,000)+\frac{1}{3}(0)+\frac{1}{3}(0)=\$10,000
\end{equation*}

To calculate the CVA use the formula: $CVA\approx LGD\cdot PD\cdot EE$:
\begin{align*}
LGD=1-R=1-0.40=0.60 \\
CVA=0.60\cdot 0.05\cdot 10,000 \\
CVA=0.03\cdot 10,000=\$300
\end{align*}
\part For an oil producer, Scenario C (MtM=-\$30,000) is when they are likely to default, but the bank's exposure is zero, so no WWR. For it to occur the counterparty should be an oil consumer (like an airline), hence Scenario A would create WWR.
\part The adjustment for the bank's own default risk is called DVA (Debit Valuation Adjustment). The total adjustment is BCVA = CVA-DVA.
\end{parts}
\end{solution}

%%%%%%%%%%%%%%%%%%%%%%%%%%%%%%%%%%%%%
\question 
A bank has a derivative position with a counterparty that has a positive Expected Exposure (EE) of \$2,000,000 over a 1-year horizon. To calculate the CVA, the bank looks at the market price of the counterparty's credit risk.

Market Data:
\begin{itemize}
\item CDS Spread ($s$): 300 basis points (3\%) per annum;
\item Recovery Rate ($R$): 40\%;
\item Risk-free rate ($r$): 0\% (to simplify discounting).
\end{itemize}

\begin{parts}
\part Using the \textit{Credit Spread Approximation}, calculate the market-implied Hazard Rate ($\lambda$) for the counterparty. 

\textbf{Hint:} $s\approx \frac{\lambda}{(1-R)}$.

\part Calculate the probability that the counterparty defaults within the 1-year period ($PD_{0,1}$) using the exponential distribution formula: $PD=1-e^{-\lambda T}$.  (If you want you can use the linear approximation for $PD$).

\part Calculate the total CVA for the position.

\part Suppose the market liquidity for the counterparty's debt dries up, and the CDS spread jumps from 300 bps to 450 bps. What is the "Credit Spread Delta" (the change in CVA for a 1 bp move in the spread) ?

\part In 2008, many banks calculated CVA using the historical PD of counterparties (based on 10 years of data) rather than the market-implied PD from CDS spreads. Explain why this led to a massive "CVA volatility" problem during the crisis.
\end{parts}
\makeemptybox{20cm}

\begin{solution}
\begin{parts}
\part Convert the spread to decimals: $s=0.03$. Identify Loss Given Default: $LGD = 1-0.40=0.60$. Finally solve for $\lambda$:
\begin{equation*}
\lambda \approx \frac{s}{(1-R)} = \frac{0.03}{0.60} = 0.05
\end{equation*}
\part Using the linear approximation: $PD\approx 0.05\cdot 1 =5\%$. The exact calculation instead gives:
\begin{equation*}
PD=1-e^{-0.05}\approx 1-0.9512=4.88\%
\end{equation*}
\part Using the linear PD for simplicity:
\begin{equation*}
CVA=LGD\cdot PD\cdot EE = 0.60\cdot 0.05\cdot 2,000,000 = \$60,000
\end{equation*}
\part Calculate the new CVA when $s$ increases by 1.5x (from 3\% to 4.5\%), and $LGD$ and $EE$ stay constant. The CVA also increases by 1.5x: $CVA_{new}=\$60,000\cdot 1.5=\$90,000$.
The total change in CVA (\$30,000) corresponds to a total change in bps of 150 bps, so
\begin{equation*}
\Delta_{spread} =\frac{30,000}{150}=\$200 \text{ per bp.}
\end{equation*}
\part Historical PDs are "backward-looking" and change very slowly. Market-implied PDs (from CDS) are "forward-looking" and react instantly to fear. During the crisis, CDS spreads skyrocketed. Banks using historical PDs were forced to suddenly switch to market prices, leading to massive, unexpected write-downs as their CVA charges jumped tenfold overnight. This is why Basel III now requires capital for CVA Volatility Risk.
\end{parts}
\end{solution}

%%%%%%%%%%%%%%%%%%%%%%%%%%%%%%%%%%%%%
\question 
An analyst wants to simulate the short term interest rate using the Vasicek model:
\begin{equation*}
dr_t = \kappa(\theta - rt)dt + \sigma dW_t
\end{equation*}

with the following parameters:
\begin{itemize}
\item asymptotic value ($\theta$): 5\%;
\item mean reversion speed ($\kappa$): 2;
\item volatility ($\sigma$): 2\%;
\item initial interest rate ($r_0$) 7\%;
\item time-horizon ($T$): 0.1 year (with 2 steps: $\Delta t= 0.05$).
\end{itemize}

Solve:
\begin{parts}
\part using the Euler scheme compute $r_1$ and $r_2$ at the two time steps. The stochastic component can be computed using the following samples from the normal distribution $Z\approx \mathcal{N}(0,1)$:
\begin{itemize}
\item Step 1 ($t=0.05$): $z_1=1.0$;
\item Step 2 ($t=0.10$): $z_2=-1.5$.
\end{itemize}

\part What would it happen if the mean reversion speed was very high (e.g. $\kappa = 50$) and keeping $\Delta t=0.05$ ?
\part Suppose to simulate an asset price using a Geometric Brownian Motion (GBM): $dS_t= \mu S_t dt+\sigma S_t dW_t$.
Explain why there is necessarily a bias in the simulation using the Euler scheme with respect to the exact solution especially for large $\Delta t$ ?
\textbf{Hint}: compare the functional forms of the GBM and of its Euler approximation.
\end{parts}
\makeemptybox{20cm}

\begin{solution}
\begin{parts}
\part The Euler formula for the Vasicek model is: $r_{t+1}=r_t +\kappa (\theta - r_t)\Delta t+\sigma \Delta t z_{t +1}$
\begin{itemize}
\item Solution step 1 ($r_1$):
\begin{gather*}
\Delta r= 2(0.05-0.07)(0.05)+(0.02)0.05(1.0) =0.00247 \\
r_1=0.07+0.00247=7.247\%
\end{gather*}
\item Solution step 2 ($r_2$):
\begin{gather*}
\Delta r = 2(0.05-0.07247)(0.05)+(0.02)0.05(-1.5) =-0.00895 \\
r_2=0.07247-0.00895=6.352\%
\end{gather*}
\end{itemize}
\part If $\kappa\Delta t >1$, the "correction" to the short rate is bigger than the distance to the asymptotic value $\theta$, forcing the solution to oscillate around the mean. %If $\kappa \Delta t >2$, Euler scheme becomes unstable and the short rate diverges. 
\part The Euler scheme is essentially a linear approximation of the precess (i.e.  in this case an exponential process). Since the exponential is a convex function, the linear approximation will tend to systematically underestimate the expected price of the asset ($\mathbb{E}[S_t]$) in the long run. To compensate for this effect is usually discretized the process $\ln(S_t)$.
\end{parts}
\end{solution}

%%%%%%%%%%%%%%%%%%%%%%%%%%%%%%%%%%%%%
	
\question 
Consider the following market data:
\begin{itemize}
\item risk free rate ($R_f$): $3\%$
\item expected return of the market ($E[R_m]$): $10\%$
\item market variance ($\sigma^2_m$): $0,04$
\end{itemize}

Using the CAPM:
\begin{parts}
\part compute the expected return of asset A, knowing its $\beta_A = 1.5$.
\part Consider asset B, which currently has a return of $8,6\%$. From a historical analysis the asset has $\beta_B = 0,9$. Determine if the asset is correctly valued, undervalued or overvalued using Jensen-alpha ($\alpha = R_{B,mkt} - E[R_B]$).
\part Decompose asset A total variance ($\sigma^2_A$) among systematic and idiosynchratic knowing that its residual variance is $\sigma^2(\epsilon_A) = 0,05$.
\part An investor builds a portfolio made of $40\%$ of asset A and $60\%$ in a risk free asset. Compute $\beta_P$ and its expected return.
\end{parts}
\makeemptybox{15cm}

\begin{solution}
\begin{parts}
\part Using the Security Market Line (SML) equation:
\begin{equation*}
\mathbb{E}[R_A] = R_f + \beta_A (E[R_m] - R_f)
\end{equation*}
Replacing the values:
\begin{equation*}
\mathbb{E}[R_A] = 0,03 + 1,5 \cdot (0,10 - 0,03) = 0,03 + 0,105 = 13,5\%
\end{equation*}
\part To valuate asset B we need to comapre its market return ($R_{B,mkt}$) with the theoretical return predicted by the CAPM ($E[R_B]$):
\begin{equation*}
\mathbb{E}[R_B] = 0,03 + 0,9 \cdot (0,10 - 0,03) = 0,03 + 0,063 = 9,3\%
\end{equation*}
Finally the Jensen-alpha:
\begin{equation*}
\alpha_B = R_{B,mkt} - E[R_B] = 8,6\% - 9,3\% = -0,7\%
\end{equation*}
Since $\alpha_B < 0$, the asset is \textbf{overvalued}. Its actual return is not enough to remunerate the associated systematic risk ($\beta=0,9$).
\part To decompose the variance in a linear model:
\begin{equation*}
\sigma^2_A = \beta_A^2 \sigma^2_m + \sigma^2(\epsilon_A)
\end{equation*}
\begin{itemize}
\item \textbf{Systematic risk}: $\beta_A^2 \sigma^2_m = (1,5)^2 \times 0,04 = 2,25 \times 0,04 = 0,09$
\item \textbf{Idiosynchratic risck}: $\sigma^2(\epsilon_A) = 0,05$
\item \textbf{Total variance}: $0,09 + 0,05 = 0,14$
\end{itemize}
\part The portfolio beta is the weighted average of asset betas:
\begin{equation*}
\beta_P = w_A \beta_A + w_{rf} \beta_{rf} = 0,4 \times 1,5 + 0,6 \times 0 = 0,6
\end{equation*}
The expected return can be computed with the SML using $\beta_P$:
\begin{equation*}
E[R_P] = R_f + \beta_P (E[R_m] - R_f) = 0,03 + 0,6 \times (0,10 - 0,03) = 7,2\%
\end{equation*}
\end{parts}
\end{solution}

%%%%%%%%%%%%%%%%%%%%%%%%%%%%%%%%%%%%%

\question 
Evaluate a Payer Interest Rate Swap (IRS) with a notional $N = €10000000$, mautirty 1 year, and annual payments. The IRS fixed rate is $K = 3\%$.
At $t=0$, the 1y-LIBOR rate and the risk free rate are both $4\%$. Determine the Swap NPV.
\makeemptybox{10cm}

\begin{solution}
A payer IRS value can be seen as the difference between a variable rate bond and a fixed rate bond:
\[ NPV = V_{float} - V_{fix} \]
			
Since the Libor rate $L$ is equal to the discounting rate $r$ the floating leg equals par:
\[ V_{float} = \frac{N(1 + L)}{1 + r} = \frac{10000000 \times 1.04}{1.04} = 10000000 \]
			
The fixed leg instead:
\[ V_{fix} = \frac{N(1 + K)}{1 + r} = \frac{10000000 \times 1.03}{1.04} = 9903846 \]
			
Hence the NPV:
\[ NPV = 10000000 - 9903846 = 96154 \]

Since interest rates increased the party paying fixed has a gain, she pays 3\% and receives 4\%.
\end{solution}

%%%%%%%%%%%%%%%%%%%%%%%%%%%%%%%%%%%%%
		
\question 
Imagine a Payer IRS with a notional $N = €10000000$, 3-year maturity and annual payments. The swap fixed rate is $K = 2.0\%$.

Due to macroeconomic tensions, the interest rate term structure steepened:
\begin{itemize}
\item $L_1$ decreases to $0.50\%$
\item $L_2$ increases to $2.50\%$
\item $L_3$ increases al $3.50\%$
\end{itemize}

Assuming constant discount factors to: $DF_1=0.98$, $DF_2=0.96$, $DF_3=0.94$:
		
\begin{parts}
\part compute the IRS swap-rate;
\part determine the swap NPV for the party paying fixed and explain why the NPV is positive even if the short-term rate decreases.
\end{parts}
\makeemptybox{10cm}
		
\begin{solution}
\begin{parts}
\part The swap-rate ($S_{mkt}$) is the weighted average of the forward rates with respect to the discount factors:
\[ S_{mkt} = \frac{\sum_{i=1}^3 DF_i \cdot L_i}{\sum_{i=1}^3 DF_i} \]
\[ S_{mkt} = \frac{(0.98 \cdot 0.005) + (0.96 \cdot 0.025) + (0.94 \cdot 0.035)}{0.98 + 0.96 + 0.94} \]
\[ S_{mkt} = \frac{0.0049 + 0.024 + 0.0329}{2.88} = \frac{0.0618}{2.88} \approx 2.146\% \]
\part The contract value results
\[ NPV = N \times (S_{mkt} - K) \times \text{Annuity} \]
\[ NPV = 10000000 \times (0.02146 - 0.0200) \times 2.88 = €42048 \]
The NPV is positive because a 3y-swap has a greater duration for longer maturities. 2 and 3 years interest rates increase overcome the decrease at 1 year, globally increasing the swap value.
\end{parts}
\end{solution}

%%%%%%%%%%%%%%%%%%%%%%%%%%%%%%%%%%%%%

\end{questions}
\end{document}
